\documentclass[main.tex]{subfiles}
\begin{document}
\section*{Appendix}
\appendix
\addcontentsline{toc}{section}{Appendix}
\setcounter{subsection}{0}
\renewcommand\thesubsection{\Alph{subsection}}
\subsection{Commitment Game Equilibria in Pure Strategies}\label{sec:proof1}
In Section~\ref{sec:GamePayoffs}, I noted that pure strategy equilibria of the base game carry over to the commitment game. In this section, I formally establish this connection by proving two supporting lemmas. Lemma~\ref{lem:SPNE_characterization} shows that any pure strategy equilibrium of the base game, such as (Straight, Straight) in the Game of Chicken, also constitutes an equilibrium outcome in the commitment game. In contrast, Lemma~\ref{lem:SPNE_characterization_2} demonstrates that strategy profiles not corresponding to any pure strategy equilibrium in the base game, such as (Move, Move), cannot arise as pure strategy equilibrium outcomes in the commitment game. Notably, this second result holds only for symmetric $2\times2$ games, and Section~\ref{sec:beyond2x2} presents a $3\times3$ game where this is violated.
\begin{lemma}[]\label{lem:SPNE_characterization}
    Consider a two player normal form game where each player has strategies $A=\{A_1,\dots,A_n\}$. Let $\{A_i, A_{-i}\}$ be a Nash equilibrium in this game. When this normal form game is played with commitment,
    \begin{enumerate}
        \item there is a strategy pair with both players committing to $A_i$ and $A_{-i}$ respectively, that forms an SPNE of the commitment game.
        
        \item if $U_i(A_i, A_{-i})\geq U_i(A'_i, B_{-i}(A_i'))\quad \forall A_i'\in A$, then there is a strategy pair where $i$ commits to $A_i$ and $-i$ waits and plays $A_{-i}$ that forms an SPNE of the commitment game.
        
        \item if $U_i(A_i, A_{-i})\geq U_i(A'_i, B_{-i}(A_i'))\quad \forall i\in\{1,2\},\; \forall A_i'\in A$, then there is a strategy pair where both players wait and then play $A_i$ and $A_{-i}$ that is an SPNE of the commitment game.
    \end{enumerate}
\end{lemma}
\paragraph{Proof:} I prove each part by showing that the constructed strategy profile yields the outcome $(A_i,A_{-i})$ and that no unilateral deviation is profitable.
 
\textit{1.} Suppose both players commit to $A_i$ and $A_{-i}$, which gives players payoffs $U(A_i,A_{-i})$. Given the other players commitment, player $i$ can unilaterally only deviate to receive any $U_i(A'_i,A_{-i})$ with $A'_i\in A$. But since $(A_i,A_{-i})$ is a Nash equilibrium of the normal‐form game, no such unilateral deviation can improve $i$'s payoff. By symmetry, this is an equilibrium. To make it subgame perfect, the strategies must also specify equilibrium strategies off-path, for which there are multiple options. All such strategy profiles are subgame perfect.
 
\textit{2.} Let player $i$ commit to $A_i$ while player $-i$ waits and plays $A_{-i}$ after observing $i$’s commitment. If $i$ deviates by committing to any $A_i'\in A$ instead, $-i$ will best respond with $B_{-i}(A_i')$ which yields $U_i\bigl(A_i',B_{-i}(A_i')\bigr)$. By assumption this does not exceed $U_i(A_i,A_{-i})$. Similarly, if $-i$ deviates by committing in period~1, the Nash equilibrium property of $(A_i,A_{-i})$ ensures she cannot gain. Thus no player has a profitable deviation.
 
\textit{3.} If both players wait and then, after the commitment period, play $A_i$ and $A_{-i}$, the outcome is again $(A_i,A_{-i})$. Any deviation to some commitment $A_i'$ would lead to the opponent responding with $B_{-i}(A_i')$, yielding a payoff $U_i\bigl(A_i',B_{-i}(A_i')\bigr)$, which is less than $ U_i(A_i,A_{-i})$ by assumption. Thus this is an SPNE: both players get the best equilibrium outcome if they wait, so committing earlier offers no advantage.
\begin{lemma}\label{lem:SPNE_characterization_2}
    Consider a two player base game, that is a symmetric $2\times2$ game with strategies $\{A_1,A_2\}$ for each player. If a strategy pair $\{A_i,A_{-i}\}$ does not form a Nash equilibrium in the base game, it will not feature in an SPNE of the game with commitment.
\end{lemma}
\paragraph{Proof:}
I prove the contrapositive. If $\{A_i,A_{-i}\}$ features in a commitment game SPNE, it must feature in a base game Nash equilibrium.
 
\textit{1.} Let's say players commit to $A_i$ and $A_{-i}$ in the commitment game SPNE. Then, by virtue of this being an SPNE it holds that $U_i(A_i,A_{-i})\geq u_i(A_i',A_{-i})\quad \forall i\in\{1,2\},\; \forall A_i'\in A$. Thus, $\{A_i,A_{-i}\}$ is a base game Nash equilibrium.
 
\textit{2.} Let's say players wait and then play $A_i$ and $A_{-i}$ in the commitment game SPNE. Then, by virtue of this being an SPNE it holds that $U_i(A_i,A_{-i})\geq U_i(A_i',A_{-i})$ $\forall i\in\{1,2\},\; \forall A_i'\in A$. Thus, $\{A_i,A_{-i}\}$ is a base game Nash equilibrium.
 
\textit{3.} This is where the $2\times2$ assumption comes in. Let's say player $i$ commits to $A_1$ and $-i$ waits. If $-i$ best responds with $A_1$,  by symmetry $\{A_1,A_1\}$ are mutual best responses and thus a Nash equilibrium of the base game. If $-i$ best responds with $A_2$, one must consider what would have happened if $i$ had committed $A_2$ instead of $A_1$. If $-i$'s best response to $A_2$ would have been $A_2$, by symmetry $i$ would have preferred to commit $A_2$ instead of $A_1$, a contradiction. Thus, $-i's$ best response to $A_2$ must be $A_1$, and thus by symmetry $\{A_1,A_2\}$ is a base game Nash equilibrium. By symmetry, this also holds for $\{A_2,A_1\}$. This concludes the proof.
 
\subsection{Proof of Proposition \ref{theo:GamePayoffs}}\label{sec:prooftheo1}
Consider the simplified $2\times2$ symmetric normal form game from table \ref{tab:general2x2}. Let $A_{1}$ denote the strategy where a player makes a commitment to $A_1$, and let $A_{2}$ denote the strategy where a player commits to $A_2$. $W$ denotes the strategy of waiting. 
To prove Proposition \ref{theo:GamePayoffs}, I will derive the mixed subgame perfect equilibria of the symmetric $2\times2$ games. In these equilibria, if a player waits, she will best respond to any commitment the other player makes. If both players wait, they will play the MSNE of the base game which yields expected payoff $\tfrac{\bparam\cparam}{\cparam+\bparam-1}$. 
Lemma \ref{lem:uniqueness} and Lemma \ref{lem:no-msne-with-pure-strats} (stated and proven further below) ensure that this is the unique mixed subgame perfect equilibrium. 
 
\paragraph{Battle of Sexes ($\bparam>1$ and $\cparam>1$)}
If player $-i$ mixes between $A_1$, $A_2$ and $W$ with respective probabilities $p$, $q$, and $1-p-q$, then player $i$'s pure strategy payoffs are
\begin{align*}
U(A_1) &= p\cdot 0 + q\cdot \bparam + (1-p-q)\cdot \bparam,\\[0mm]
U(A_2) &= p\cdot \cparam + q\cdot 1 + (1-p-q)\cdot \cparam,\\[0mm]
U(W) &= p\cdot \cparam + q\cdot \bparam + (1-p-q)\cdot \tfrac{\bparam\cparam}{\bparam+\cparam-1}.
\end{align*}
For player $i$ to mix over these strategies, she must be indifferent between them. The $p$ and $q$ that equate the pure strategy payoffs are
\begin{alignat*}{1}
    p &= \frac{\bparam(1-\bparam)^2}{(\bparam+\cparam-1)(\bparam^2-\bparam+\cparam^2-\cparam)}\\[6pt]
    q &= \frac{c-(1-p)b}{\cparam-1}
\end{alignat*}
It can be shown that $p,q\in[0,1]$. By symmetry, this holds for player $-i$ as well. To obtain the equilibrium payoff, plug $p$ and $q$ into any of the above pure strategy payoffs. Consider $U(W)$. Note that $\bparam$ and $\cparam$ are both bigger than $\tfrac{\bparam\cparam}{\cparam+\bparam-1}$. Thus, the payoff in the commitment game equilibrium exceeds the payoff in the base game.
\paragraph{Game of Chicken ($\bparam>1$ and $1\geq \cparam>0$)}
For $1\geq\cparam$, waiting weakly dominates $A_2$. Thus, in the mixed subgame perfect equilibrium there is mixing only over $A_1$ and $W$. The probability $p$ that equates $U(A_1)$ and $U(W)$ is 
\[
p=\frac{\bparam^2-\bparam}{\bparam^2-\bparam+\cparam^2-\cparam+\bparam\cparam}.
\]
It can be shown that $p\in[0,1]$. Consider $U(W)$. Since $\cparam\leq1\Leftrightarrow \cparam\leq\tfrac{\bparam\cparam}{\cparam+\bparam-1}$. Thus, the commitment game payoff is worse than the base game payoff for $\cparam<1$, except for the edge case with $c=1$ where it is the same.
\vspace{0.5cm}
\newline
This concludes the proof of Proposition~\ref{theo:GamePayoffs}. The next equilibria I derive prove Corollary~\ref{cor:staghunt-coord}.
\paragraph{Coordination Game ($\cparam<0$ and $\bparam<0$)}
In the Coordination Game and the Stag Hunt the best responses to commitments are different from the best responses in the Game of Chicken and the Battle of Sexes. Here, if player $-i$ mixes over the strategies $A_1$, $A_2$, and $W$ with probability $p$, $q$, and $1-p-q$, player $i$'s pure strategy payoffs are
\begin{align*}
U(A_1) &= p\cdot 0 + q\cdot \bparam + (1-p-q)\cdot 0,\\[0mm]
U(A_2) &= p\cdot \cparam + q\cdot 1 + (1-p-q)\cdot 1,\\[0mm]
U(W) &= p\cdot 0 + q\cdot 1 + (1-p-q)\cdot \tfrac{\bparam\cparam}{\bparam+\cparam-1}.
\end{align*}
In a mixed equilibrium with mixing over these three strategies, all pure strategy payoffs must be equal. The $p$ and $q$ that equate the above pure strategy payoffs are
\begin{alignat*}{1}
p &= \frac{(1-\bparam)\,(bc-b-c+1)}{(1-\bparam-\cparam)(2\bparam\cparam-\bparam - \cparam + 1)},\\[1mm]
q &= \frac{p(1-c)-1}{-b}\,.
\end{alignat*}
It can be shown that $p,q\in[0,1]$. To compare the equilibrium payoff of the commitment game to the base game MSNE payoff, plug $p$ and $q$ into $U(W)$. Since $0>\tfrac{\bparam\cparam}{\cparam+\bparam-1}$, the commitment game payoff exceeds the MSNE base game payoff.
 
However, note that committing to $(A_1,A_1)$ or players waiting and then playing ($A_1,A_1$) also form symmetric equilibria in the commitment game. These equilibria yield payoff $1$ for each player, which is bigger than $U(W)$.
\paragraph{Stag Hunt ($0\leq\bparam<1$ and $\cparam<0$)}
Here, committing to $A_1$ is weakly dominated by waiting. $A_1$ will not feature in a mixed equilibrium and $p$ can be set to $0$. But then, $A_2$ weakly dominates waiting. Thus, there is no equilibrium in mixed strategies when commitment is introduced in a Stag Hunt game.
\paragraph{Dominant Strategies}
When $A_1$ or $A_2$ is a (weakly) dominant strategy in the game without commitment, there is little to explore. In equilibrium, players can commit to their dominant strategy, they can wait before playing it or they can mix between committing to it and waiting and then playing it. These equilibria do not offer interesting strategic dynamics, but are mentioned for completeness.
\subsection{Proof of Proposition \ref{theo:convergence}}\label{sec:prooftheo2}
This section proves Proposition~\ref{theo:convergence}. I prove 3 lemmas. 
Lemma \ref{lem:uniqueness} states that the proposed mixed equilibrium is the unique subgame perfect equilibrium where players mix over committing and waiting, assuming that players play the base game MSNE if both wait. 
Lemma \ref{lem:no-msne-with-pure-strats} shows that when players wait and then play a base game equilibrium in pure strategies, no mixed equilibrium exists in which both players mix over commitment and waiting. 
Lemma \ref{lem:convergence} then describes the unique mixed equilibrium as $T\to\infty$ and characterizes closed form expressions for all $p_t$, $q_t$ and $\sigma(W)$ in terms of $\bparam,\cparam$ and $T$.

Note that I refer to uniqueness in terms of being the only equilibrium where both players mix over committing and waiting. This definition is more restrictive than necessary. I conjecture that the equilibrium described in Proposition 2 is the unique equilibrium in which miscoordination can occur, which means that both players assign positive probability to committing in some period t. 
Other equilibria can exist in the game, but they are either in pure strategies or mixed in a trivial sense. For instance, in an equilibrium where Player 1 mixes over committing in period 1 and 2, while Player 2 mixes over committing in period 3 and waiting, Player 1 effectively induces her preferred outcome like in a pure strategy equilibrium.


\begin{lemma}[Unique Interior MSNE]\label{lem:uniqueness}
Consider the general symmetric $2\times2$ game in~\autoref{tab:general2x2}, with parameters $\bparam > 1$ and $\cparam > 0$. Let each player have the ability to commit to one of two actions, $A_1$ or $A_2$, in any of $T$ commitment periods, or to wait until the end. In the subgame where both players wait, they play the mixed-strategy Nash equilibrium of the base game. Then, the commitment game has a unique mixed subgame perfect Nash equilibrium where players mix over committing and waiting. Specifically:
\begin{enumerate}
    \item If $\cparam > 1$, both players mix over all strategies $A_{1t}$, $A_{2t}$, and $W$.
    \item If $\cparam \leq 1$, both players mix over all strategies $A_{2t}$ and $W$.
\end{enumerate}
\end{lemma}
\paragraph{Proof:}
Remember that $U(A_{1t})$ denotes the payoff of committing to strategy $A_1$ in period $t$. $p_t$ and $q_t$ denote the probability of a player playing strategy $A_{1t}$ or $A_{2t}$, respectively, and $\sigma(W)=1-\sum_{\tau=1}^T(p_\tau+q_\tau)$ denotes the probability of waiting. Then, the pure strategy payoffs in a commitment game with $T$ periods are
\begin{align}
U(A_{1t}) &= \sum_{\tau=1}^{t-1}\left( p_\tau \cparam + q_\tau \bparam \right) + p_t \cdot 0 + q_t \cdot \bparam + \left(1 - \sum_{\tau=1}^{t}(p_\tau + q_\tau)\right) \bparam
\label{eq:U_A1t} \\[6pt]
U(A_{2t}) &= \sum_{\tau=1}^{t-1}\left( p_\tau \cparam + q_\tau \bparam \right) + p_t \cparam + q_t \cdot 1 + \left(1 - \sum_{\tau=1}^{t}(p_\tau + q_\tau) \right) \cparam
\label{eq:U_A2t} \\[6pt]
U(W) &= \sum_{\tau=1}^{T} \left( p_\tau \cparam + q_\tau \bparam \right) + \sigma(W) \left( \frac{\bparam\cparam}{\bparam + \cparam - 1} \right)
\label{eq:U_W}
\end{align}
To illustrate, suppose player 1 plays $A_{1t}$. With probability $\sum_{\tau=1}^{t-1}(p_\tau + q_\tau)$, player 2 commits to $A_1$ or $A_2$ before her, and player 1 best responds, receiving $\cparam$ or $\bparam$ accordingly. 
\footnote{I adopt the convention that $\sum_{\tau=1}^{t-1} (\cdots) = 0$ when $t=1$, i.e.\ the sum is empty and thus vanishes.}
With probability $p_t + q_t$, both players commit simultaneously: if 2 also commits $A_1$, the payoff is $0$; if 2 commits $A_2$, player 1 receives $\bparam$. With the remaining probability $1 - \sum_{\tau=1}^{t}(p_\tau + q_\tau)$, player 2 was planning to move only after player 1 committed $A_1$. In this case, player 1 commits, player 2 best responds with $A_2$ and 1 receives the payoff $\bparam$.
 
For now, consider the case $\cparam > 1$.
I am looking for equilibria where players mix over committing and waiting. Thus, $\sigma(W)$ and some $p_t$ or $q_t$ must be greater than 0.
First, consider
\begin{align*}
U(A_{1, T}) - U(W) &= p_{T} \cdot (- \cparam) + \sigma(W) \cdot \left( \frac{\bparam(\bparam-1)}{\bparam+\cparam-1} \right) \\
U(A_{2, T}) - U(W) &= q_{T} \cdot (1-\bparam) + \sigma(W) \cdot \left( \frac{\cparam(\cparam-1)}{\bparam+\cparam-1} \right)
\end{align*}
If $p_{T}=0$ and $\sigma(W)>0$, $U(A_{1,T})>U(W)$. It must be that players play $A_{1T}$ with positive probability, otherwise waiting would be dominated. The same holds for $A_{2T}$, thus $p_T>0$ and $q_T>0$.
Next, consider
\begin{align*}
    U(A_{1t})-U(A_{1,t+1}) &= p_t \cdot (-\cparam) + p_{t+1} \cdot \bparam \\
    U(A_{2t})-U(A_{2,t+1}) &=  q_t \cdot (1-\bparam) + q_{t+1} \cdot (\cparam-1)
\end{align*}
If $p_t=0$ and $p_{t+1}>0$, $U(A_{1t})>U(A_{1t+1})$. Thus for a player to mix over both $A_{1t}$ and $A_{1t+1}$, it must be that the other player mixes over both $A_{1t}$ and $A_{1t+1}$ with positive probability. Since both players must have probability $p_T>0$ as argued above, they must also have probability $p_{T-1}$ greater than 0. It then follows that they must have all probabilities $\{p_t\}_{t=1}^T$ greater than 0. The same holds for $A_{2t}$ and $q_t$.
Thus, the only equilibrium where both players mix over committing and waiting is such that both players mix over all strategies $\{A_{1t}, A_{2t}, W\}_{t=1}^{T}$ with positive probability. This concludes the proof for $\cparam>1$.

When $\cparam\leq1$, playing any $A_{2t}$ is weakly dominated by waiting: equations \ref{eq:U_A2t} and \ref{eq:U_W} imply $U(A_{2t})\leq U(W)$. 
Players will not put probability weight on $A_2$ and thus all $q_t=0$. The argument of the proof for $\cparam>1$ then holds for $A_1$ and $W$. In equilibrium, players will mix across all strategies in the set $\{A_{1t}, W\}_{t=1}^{T}$. This concludes the proof.

\begin{lemma}\label{lem:no-msne-with-pure-strats}
    Consider the general symmetric $2\times2$ game in~\autoref{tab:general2x2}, with parameters $\bparam > 1$ and $\cparam > 0$. 
    Let each player have the ability to commit to one of two actions, $A_1$ or $A_2$, in any of $T$ commitment periods, or to wait until the end. 
    In the subgame where both players wait, they play one of the pure strategy Nash equilibria of the base game.
    Then, as long as $\bparam\neq\cparam$, the commitment game does not have a mixed subgame perfect Nash equilibrium where both players mix over committing and waiting.
\end{lemma}
\paragraph{Proof:}
I aim to show that there is no mixed equilibrium in which both players randomize between committing and waiting. 
In any mixed equilibrium, players must be indifferent between all strategies they assign positive probability to.
I will show that no such indifference condition can be satisfied when both player mix between commitment and waiting strategies.

Let each player mix over strategies $A_{1t}$, $A_{2t}$, and $W$ with probabilities $p_t^i$, $q_t^i$, and $\sigma(W)^i$ for player $i$, and analogously for player $-i$. Then it holds that
\begin{align}
    U_i(A_{1, t}) - U_i(W) 
    &= \sum_{\tau = t+1}^T p_{\tau}^{-i}(b-c) + p_{t}^{-i}(-c) + \sigma(W)^{-i} \bigl(b - \max\{b,c\}\bigr),
    \label{eq:Udiff_A1t} \\[6pt]
    U_i(A_{2, t}) - U_i(W) 
    &= \sum_{\tau = t+1}^T q_{\tau}^{-i}(c-b) + q_{t}^{-i}(1-b) + \sigma(W)^{-i} \bigl(c - \max\{b,c\}\bigr),
    \label{eq:Udiff_A2t}
\end{align}
and analogously for player $-i$:
\begin{align}
    U_{-i}(A_{1, t}) - U_{-i}(W) 
    &= \sum_{\tau = t+1}^T p_{\tau}^{i}(b-c) + p_{t}^{i}(-c) + \sigma(W)^{i} \bigl(b - \min\{b,c\}\bigr),
    \label{eq:Udiff_A1t_-i} \\[6pt]
    U_{-i}(A_{2, t}) - U_{-i}(W) 
    &= \sum_{\tau = t+1}^T q_{\tau}^{i}(c-b) + q_{t}^{i}(1-b) + \sigma(W)^{i} \bigl(c - \min\{b,c\}\bigr).
    \label{eq:Udiff_A2t_-i}
\end{align}
Assume $\bparam > \cparam$. Then from equations~\eqref{eq:Udiff_A2t} and~\eqref{eq:Udiff_A2t_-i}, it follows:
\[
U_i(A_{2t}) < U_i(W) \quad \text{and} \quad U_{-i}(A_{2t}) \leq U_{-i}(W) \quad \text{for all } t.
\]
Thus, neither player will commit to any $A_{2t}$ with positive probability in any mixed equilibrium. The case $\bparam < \cparam,$ is symmetric and implies that neither player will commit to $A_1$. For the remainder of the proof, I focus on the case $b > c$.

Suppose for contradiction that one player mixes over some, but not all, commitment strategies. Then there exists some $t$ such that $p_t = 0$ for one player. Consider the difference in expected utility between adjacent commitment periods:
\begin{equation}
    U(A_{1t})-U(A_{1,t+1}) = p_t \cdot (-\cparam) + p_{t+1} \cdot \bparam \label{eq:A1-difference-lemma4}
\end{equation}
Since $p_t = 0$ but $p_{t+1} > 0$, by iterating equation \ref{eq:A1-difference-lemma4} it follows that for the other player $U(A_{1t})$ is bigger than all $U(A_{1t+1}),\dots,U(A_{1T})$. By equation \ref{eq:Udiff_A1t} it follows that $U(A_{1t})>U(W)$. Therefore, $A_{1t}$ dominates commitment after period $t$ and waiting, contradicting indifference. Thus, for $\sigma(W) > 0$ to hold, both players must assign positive probability to all commitment strategies.
However, consider 
\begin{align*}
    &U_i(A_{1, T}) - U_i(W) = p_{T}^{-i} \cdot (- \cparam) + \sigma(W)^{-i} \cdot (b-\max\{b,c\})
\end{align*}
Since $b > c$, the second term is zero. If player $-i$ assigns positive probability to $A_{1T}$, i.e., $p_T^{-i} > 0$, then $U_i(A_{1,T}) < U_i(W)$. This implies that player $i$ cannot be indifferent between waiting and $A_{1T}$ and thus will prefer waiting to committing in any commitment period.

Therefore, it is not possible that both players mix over all commitment strategies and waiting in equilibrium. This concludes the proof.

\begin{lemma}\label{lem:convergence}
    Consider the general symmetric $2\times2$ game in~\autoref{tab:general2x2} with parameters satisfying $\bparam>1$ and $\cparam>0$. The mixed subgame perfect Nash equilibrium corresponds to the strategy profile described in Table~\ref{tab:lemma4-probabilities}. Table~\ref{tab:convergence} summarizes the convergence of equilibrium probabilities and payoffs as $T\to\infty$.
\end{lemma}
\begin{table}[ht]
\centering
\begin{tabular}{llcc}
\toprule
\textbf{Strategy} & \textbf{Expression} & \multicolumn{2}{c}{\textbf{Case}} \\
& & $\cparam > 1$ & $0 < \cparam \leq 1$ \\
\midrule
Commit $A_1$ immediately & 
$p_1$ &
$\tfrac{\cparam - \bparam}{\cparam\ \tfrac{\cparam-1}{\bparam-1} \left( \tfrac{\cparam(\cparam - 1)}{\bparam(\bparam - 1)} \right)^T\ -\ b }$ &
$\tfrac{1}{\tfrac{1 - \rho^{T}}{1 - \rho} + \rho^{T}\,\tfrac{\bparam + \cparam - 1}{\bparam - 1}}$ \\
[4ex]
 
Commit $A_2$ immediately & 
$q_1$ &
$\tfrac{p_1 \bparam + \cparam - \bparam}{\cparam - 1}$ &
$0$ \\
[3ex]
 
Commit $A_1$ in period $t$ & 
$p_t$ &
$p_1 \left( \tfrac{\cparam}{\bparam} \right)^{t - 1}$ &
$p_1 \left( \tfrac{\cparam}{\bparam} \right)^{t - 1}$ \\
[3ex]
 
Commit $A_2$ in period $t$ & 
$q_t$ &
$q_1 \left( \tfrac{\bparam - 1}{\cparam - 1} \right)^{t - 1}$ &
$0$ \\
[3ex]
 
Wait & 
$\sigma(W)$ &
$\tfrac{\bparam + \cparam - 1}{\bparam-1} \left( \tfrac{\cparam}{\bparam} \right)^T p_1$ &
$\tfrac{\bparam + \cparam - 1}{\bparam-1} \left( \tfrac{\cparam}{\bparam} \right)^T p_1$ \\
\bottomrule
\end{tabular}
\caption{Equilibrium Probabilities in the Commitment Game}
\label{tab:lemma4-probabilities}
\end{table}
\paragraph{Proof}
I will first consider $\cparam>1$ and then $\cparam\leq1$.
 
\textbf{Case 1: $\cparam>1$.}
Suppose player $-i$ mixes over $A_{1t}$, $A_{2t}$ and $W$ with probabilities $p_{t}$, $q_{t}$ and $\sigma(W)=1-\sum_{\tau=1}^{T}(p_\tau+q_\tau)$. Then again, player $i$'s pure strategy payoffs are described by equations \ref{eq:U_A1t}, \ref{eq:U_A2t} and \ref{eq:U_W}.
In a mixed equilibrium, all these strategies must yield the same expected payoff:
\[
U(A_{11})=...=U(A_{1,T})=U(A_{2,1})=...=U(A_{2,T})=U(W).
\]
It can be shown that 
\begin{align}
U(A_{11}) &= U(A_{21}) 
&&\Longleftrightarrow \quad q_1 = \frac{p_1 \bparam +\cparam-\bparam}{\cparam-1}
\label{eq:initial} \\
%
U(A_{1t}) &= U(A_{1t+1}) 
&&\Longleftrightarrow \quad p_{t+1} = p_t\cdot\frac{c}{b}
\label{eq:action1} \\
%
U(A_{2t}) &= U(A_{2t+1}) 
&&\Longleftrightarrow \quad q_{t+1} = q_t\cdot\frac{b-1}{c-1}
\label{eq:action2} \\
%
U(A_{1T}) &= U(W) 
&&\Longleftrightarrow \quad \sigma(W) = p_T\cdot\frac{c}{b}\cdot\frac{b+c-1}{b-1}
\label{eq:terminal1} \\
%
U(A_{2T}) &= U(W) 
&&\Longleftrightarrow \quad \sigma(W) = q_T \cdot \frac{b-1}{c-1}\frac{b+c-1}{c}
\label{eq:terminal2}
\end{align}
I now want to express $p_t$, $q_t$ and $\sigma(W)$ as closed form solutions in terms of $\bparam$, $\cparam$ and $T$. First note that \ref{eq:action1} and \ref{eq:action2} can be used to write
\[
p_t = p_1 \left(\frac{\cparam}{\,\bparam}\right)^{t-1} \qquad 
q_t = q_1 \left(\frac{\bparam-1}{\,\cparam-1}\right)^{t-1}
\]
Next, I will use these to plug into \ref{eq:terminal1} and \ref{eq:terminal2} to obtain
\[
p_1 \, \left(\frac{\cparam}{\,\bparam}\right)^{T} \, \frac{\bparam+\cparam-1}{\bparam-1} =  q_1 \, \frac{\bparam+\cparam-1}{\cparam} \, \left(\frac{\bparam-1}{\,\cparam-1}\right)^{T}
\]
Next, I plug in \ref{eq:initial} for $q_1$ to obtain an expression that can be rearranged for $p_1$. This finally yields:
\begin{equation}
   p_1 = \frac{\cparam-\bparam}{\,\cparam\frac{c-1}{b-1}\left(\frac{\cparam(\cparam-1)}{\,\bparam(\bparam-1)}\right)^{T}-b}, \qquad
   \label{eq:p_deriv}
\end{equation}
I have thus derived a closed form expression for $p_1$, and using equations \ref{eq:initial}-\ref{eq:terminal2} I can obtain closed form expressions for all $p_t$, $q_t$ and $\sigma(W)$. Next, let's analyze the equilibrium as $T\to\infty$. Define $r=\frac{\cparam(\cparam-1)}{\bparam(\bparam-1)}$. Then
\begin{equation*}
p_1 = \frac{\cparam - \bparam}{\cparam \frac{\cparam-1}{\bparam-1}r^T-\bparam} 
 \text{,} \quad 
q_1 = \frac{p_1 \bparam + \cparam - \bparam}{\cparam - 1}
\quad \text{and} \quad 
\sigma(W) = p_1\left(\frac{\cparam}{\bparam}\right)^T\frac{\bparam + \cparam - 1}{\bparam-1}  \,.
\end{equation*}
As $T\to\infty$, I distinguish three cases:
\begin{itemize}
    \item If $\cparam<\bparam$, then $r<1$, and
    \[
    p_1 \to \frac{\bparam-\cparam}{\,\bparam},\qquad q_1 \to 0,\qquad \sigma(W) \to 0.
    \]
    \item If $\cparam>\bparam$, then $r>1$, and
    \[
    p_1 \to 0,\qquad q_1 \to \frac{\cparam-\bparam}{\,\cparam-1},\qquad \sigma(W) \to 0.
    \]
    \item Finally, if $\cparam=\bparam$
    \[
    p_1=\frac{1}{T+1}\frac{b-1}{2b-1}, \quad q_1 = \frac{1}{T+1}\frac{b}{2b-1}, \quad \sigma(W)\to 0
\]
\end{itemize}
Note that the case $b=c$ is not derived from equation \ref{eq:p_deriv}. Instead, I plug equations \ref{eq:initial}-\ref{eq:terminal2} into $\sum_{t=1}^T(p_t+q_t)+\sigma(W)=1$ to obtain closed form expressions. To see the equilibrium payoff in all cases consider $U(W)$. As can be seen, it converges to $\min\{\bparam,\cparam\}$ as $T\to\infty$.
\newline
\newline
\textbf{Case 2: $\cparam\leq1$.}
Here, committing to $A_2$ is weakly dominated by waiting. Thus all $q_t=0$ and the mixed equilibrium is only characterized by Equation \ref{eq:action1} and \ref{eq:terminal1}.
Let $\rho=\frac{\cparam}{\bparam}$. Since the probability sequence must satisfy $ \sum_{\tau=1}^{T} p_\tau + \sigma(W) = 1 $, I can write 
$$
p_1\sum_{\tau=0}^{T-1} \rho^\tau + p_1\,\rho^{T}\,\frac{\bparam+\cparam-1}{\bparam-1} = 1.
$$
The sum of the geometric series is $\sum_{\tau=0}^{T-1} \rho^\tau = \frac{1-\rho^{T}}{1-\rho}$. Plugging this into the above expression and solving for $p_1$ gives 
\[
p_1 = \frac{1}{\frac{1-\rho^{T}}{1-\rho} + \rho^{T}\,\frac{\bparam+\cparam-1}{\bparam-1}},
\]
which together with \ref{eq:action1} and \ref{eq:terminal1} give the closed form expressions for all $p_t$ and $\sigma(W)$.
Again, let me analyze the equilibrium as $T\to\infty$. Note that $\rho^{T}\to 0$, so that the geometric series simplifies to $\frac{1-\rho^{T}}{1-\rho}\to \frac{1}{1-\rho}$. Thus,
\[
p_1\to 1-\rho,\qquad \sigma(W) \to 0.
\]
Consider $U(W)$ again to see that the equilibrium payoff converges to $\cparam$. Thus fur $\cparam<1$ the equilibrium payoff worsens as $T$ increases. Note that $\cparam=1$ is an edge case, where the MSNE base game payoff $\frac{\bparam\cparam}{\bparam+\cparam-1}=1$ and thus the commitment game equilibrium payoff stays $1$ as $T\to\infty$. This concludes the proof of Lemma \ref{lem:convergence}.

\end{document}